% ============================================================
%  UNCONSTRAINED VIRTUAL-SUBSTATE PARTITION-DEPTH MINIMISATION
%  A Validation Blueprint for Force-Free, Simulation-Free
%  Single-Transient Mass Spectrometry
% ============================================================

\documentclass[twocolumn,10pt,a4paper]{article}

\usepackage[utf8]{inputenc}
\usepackage[T1]{fontenc}
\usepackage{lmodern}
\usepackage[a4paper,top=2cm,bottom=2.2cm,left=1.8cm,right=1.8cm,columnsep=0.6cm]{geometry}
\usepackage{amsmath,amssymb,amsthm,mathtools}
\usepackage{bm}
\usepackage{booktabs}
\usepackage{array}
\usepackage{enumitem}
\usepackage{graphicx}
\usepackage{xcolor}
\usepackage{microtype}
\usepackage[round,sort&compress]{natbib}
\usepackage[colorlinks=true,linkcolor=blue!60!black,
            citecolor=blue!60!black,urlcolor=blue!60!black]{hyperref}
\usepackage{fancyhdr}
\usepackage{algorithm}
\usepackage{algpseudocode}
\usepackage{float}

\pagestyle{fancy}
\fancyhf{}
\fancyhead[LE,RO]{\thepage}
\fancyhead[RE]{\small Unconstrained Virtual-Substate Depth Minimisation}
\fancyhead[LO]{\small Sachikonye}
\renewcommand{\headrulewidth}{0.4pt}

\newtheoremstyle{thmstyle}{\topsep}{\topsep}{\itshape}{}{\bfseries}{.}{.5em}{}
\theoremstyle{thmstyle}
\newtheorem{theorem}{Theorem}[section]
\newtheorem{lemma}[theorem]{Lemma}
\newtheorem{proposition}[theorem]{Proposition}
\newtheorem{corollary}[theorem]{Corollary}
\theoremstyle{definition}
\newtheorem{definition}[theorem]{Definition}
\newtheorem{axiom}{Axiom}
\newtheorem{protocol}[theorem]{Protocol}
\newtheorem{experiment}[theorem]{Experiment}
\theoremstyle{remark}
\newtheorem{remark}[theorem]{Remark}

\newcommand{\R}{\mathbb{R}}
\newcommand{\N}{\mathbb{N}}
\newcommand{\Zp}{\mathbb{Z}^{+}}
\newcommand{\Q}{\mathbb{Q}}
\newcommand{\kB}{k_{\mathrm{B}}}
\newcommand{\Parts}{\mathcal{P}}
\newcommand{\PhiMap}{\varPhi}
\newcommand{\Nstate}{N_{\mathrm{state}}}
\newcommand{\Mmax}{M_{\max}}
\newcommand{\Sspace}{\mathcal{S}}
\newcommand{\Depth}{\mathcal{M}}
\newcommand{\Lagr}{\mathcal{L}}
\newcommand{\Apartition}{\mathbf{A}_{\Depth}}
\newcommand{\partinert}{\mu}
\newcommand{\taup}{\tau_{p}}
\newcommand{\calF}{\mathcal{F}}

\title{\textbf{Unconstrained Virtual-Substate Partition-Depth Minimisation:\\
A Self-Contained Validation Blueprint for Force-Free,\\
Single-Transient Mass Spectrometry}}

\author{
Kundai Farai Sachikonye\\
Technical University of Munich\\
Chair of Proteomics and Bioanalytics\\
\texttt{kundai.sachikonye@wzw.tum.de}
}

\date{\today}

\begin{document}
\maketitle

\begin{abstract}
\noindent
We derive a complete, self-contained validation programme for the
proposition that a mass spectrometer is a partition-depth observatory
whose output is determined by a deterministic computer programme over
ternary partition addresses, rather than a force-driven instrument
requiring numerical integration of Newtonian dynamics. The derivation
proceeds from a single geometric axiom---the Bounded Phase Space
Law---and produces every result that the validation programme tests.
We prove: (i) oscillatory necessity for persistent bounded systems;
(ii) the partition coordinate system $(n,\ell,m,s)$ with shell capacity
$C(n)=2n^{2}$; (iii) the bijection $\PhiMap:\Zp\to\Parts$ between the
positive integers and the partition states, with explicit
inversion formula; (iv) the time-count identity
$dM/dt=\omega/(2\pi)$ as a definition of derived time rather than a
circular identity; (v) the triple equivalence
$\mathbf{Osc}\simeq\mathbf{Cat}\simeq\mathbf{Part}$ with explicit
conversion maps; (vi) the partition Lagrangian
$\Lagr_{\Depth}=\tfrac{1}{2}\partinert|\dot{\mathbf{x}}|^{2}+\partinert\dot{\mathbf{x}}\cdot\Apartition-\Depth$
and the four analyser scaling laws as topology specialisations of one
equation; (vii) the phase-coherence relation
$\Delta\theta=2\pi\Delta M$ between a precursor and every member of its
forward sequence $\calF(M_{\mathrm{prec}})$, with fragment subharmonics
at $\omega_{f}=\omega_{\mathrm{prec}}\sqrt{m_{\mathrm{prec}}/m_{f}}$;
(viii) the virtual substate construction with mean-recovery constraint
$\bar{s}=s$ admitting off-shell components $s_{j}\notin[0,1]$; (ix) the
unconstrained subtask theorem ensuring local virtual values are free
under preserved global value; (x) the backward-completion algorithm in
$\mathcal{O}(\log_{3}N)$ steps with $\Theta(N)$ collapse when virtual
components are restricted to the physical cube; (xi) the six-pass
shader architecture as a complete force-free instrument. From these
derivations we extract five falsifiable experiments, each with a
predicted-versus-measured identity, a JSON record schema, a pass
criterion at the precision the framework predicts, and a clear failure
mode. The blueprint is implementable on standard \texttt{mzML} and raw
Orbitrap transient data without wet-lab work. Every assertion is
proved; nothing is cited from prior work of the author; the
literature references are restricted to canonical mathematical and
physical results used as lemmas.

\medskip
\noindent\textbf{Keywords:} bounded phase space, partition bijection,
time-count identity, partition Lagrangian, phase coherence, virtual
substates, mean recovery, off-shell decomposition, backward trajectory
completion, six-pass shader pipeline, falsifiable validation,
single-transient mass spectrometry.
\end{abstract}

\tableofcontents
\bigskip\hrule\bigskip

% =========================================================
\section{Introduction}
\label{sec:intro}
% =========================================================

\subsection{Two pictures of a mass spectrometer}

A mass spectrometer is conventionally described as a force-driven
instrument. An ion source applies an electric field which strips
electrons; ion optics steer the resulting charged species through
electromagnetic potentials; a mass analyser separates ions through the
Lorentz force $\mathbf{F}=q(\mathbf{E}+\mathbf{v}\times\mathbf{B})$
\citep{griffiths2017}; a detector registers mechanical impacts.
Forward simulation packages integrate Newton's equations of motion
through field geometries to predict observables \citep{dahl2000}. A
spectrum is the result of a long causal chain of force-mediated
events; a measurement is property revelation.

This paper develops, derives, and prepares for falsification a
different picture. A mass spectrometer is a \emph{partition-depth
observatory}: ions are bounded oscillatory systems whose categorical
identity is encoded in a small ternary address, and the instrument
reads that address by completing a trajectory in a bounded
three-dimensional entropy coordinate space
$\Sspace=[0,1]^{3}$. There are no forces involved at any stage; the
``forces'' of the conventional picture are the Euler--Lagrange
equations for partition depth minimisation. The instrument is a
deterministic computer programme whose output is read from the
partition structure, not computed from dynamics.

\subsection{The validation question}

These two pictures are observationally equivalent in their predicted
$m/z$ scaling laws, but they make sharply different predictions about
what information is present in a single measurement. The force picture
treats each scan as one piece of information: an MS\textsuperscript{1}
spectrum gives $m/z$; an MS\textsuperscript{2} spectrum gives fragment
masses; a charge envelope gives $z$ states; complete characterisation
requires hours of instrument time across multiple experiments. The
partition picture predicts that a \emph{single} Orbitrap transient of
one precursor ion already contains the full MS\textsuperscript{2}
fragment spectrum (as subharmonics at irrational frequency multiples
with phase $\Delta\theta=2\pi\Delta M$), all charge state spectra (at
$\sqrt{z}$ multiples), both polarity modes (as standing-wave phase
signatures), and the complete partition history of the ion.

The validation question is sharp: are those predicted subharmonic
components present at the predicted positions with the predicted
phases in real Orbitrap transient data?

\subsection{Goals of the blueprint}

We pursue four goals in this paper.
\begin{enumerate}[leftmargin=*,itemsep=0.2em]
\item Prove from scratch every claim the partition picture makes
about an ion. The reader does not need to take any prior work on
trust; if a result is used in the validation, it is derived here.
\item Establish the bridge from the partition framework to a
computational programme. We show that the same ternary address
encodes both the ion's position and its trajectory, and that the GPU
fragment shader evaluating partition depth at cell coordinates is
performing the physical observation rather than simulating it.
\item Identify five falsifiable experiments. Each experiment maps one
derived theorem to one observable computable from \texttt{mzML} or raw
transient data, with a precision criterion the framework predicts and
a clear failure mode if the framework is wrong.
\item Specify each experiment to the level of detail needed for
implementation: input data contract, predicted statistic, measured
statistic, JSON record schema, and pass/fail rule.
\end{enumerate}

\subsection{Outline}

Section~\ref{sec:axiom} establishes the single axiom and derives
oscillatory necessity. Section~\ref{sec:partition} derives the
partition coordinate system $(n,\ell,m,s)$ and the capacity formula
$C(n)=2n^{2}$. Section~\ref{sec:bijection} proves the bijection
$\PhiMap:\Zp\to\Parts$ and the traversal/existence distinction.
Section~\ref{sec:timecount} establishes the time-count identity as a
non-circular definition of derived time.
Section~\ref{sec:triple} proves the triple equivalence and the
observation--computation identity.
Section~\ref{sec:lagrangian} derives the partition Lagrangian and the
four analyser scaling laws.
Section~\ref{sec:phase} proves phase coherence and the
subharmonic-encoding theorem.
Section~\ref{sec:virtual} develops virtual substates and the
mean-recovery constraint.
Section~\ref{sec:tensor} stacks four dimensions of virtual
decomposition into the partition tensor.
Section~\ref{sec:backward} derives backward trajectory completion in
$\mathcal{O}(\log_{3}N)$ and the $\Theta(N)$ collapse without virtual
components.
Section~\ref{sec:shader} describes the six-pass shader pipeline as a
complete force-free instrument.
Section~\ref{sec:experiments} states the five falsifiable experiments
in protocol form.
Section~\ref{sec:contracts} specifies data contracts, JSON record
schemas, and pass criteria.
Section~\ref{sec:discussion} discusses limitations and scope.

% =========================================================
\section{The Single Axiom and Oscillatory Necessity}
\label{sec:axiom}
% =========================================================

\subsection{Phase space and the boundedness axiom}

Consider a classical dynamical system with $d$ spatial degrees of
freedom, on phase space $\Gamma\subseteq\R^{2d}$ with canonical
coordinates $(\mathbf{q},\mathbf{p})$. A measure-preserving flow
$\phi_{t}:\Gamma\to\Gamma$ generated by a Hamiltonian preserves the
Liouville measure
$\mu=d^{d}q\,d^{d}p/h^{d}$, where $h$ is a fixed phase-space cell size
fixing the units.

\begin{axiom}[Bounded Phase Space Law]
\label{ax:bps}
Every persistent dynamical system occupies a bounded region
$\Omega\subseteq\Gamma$ of finite Liouville measure
$\mu(\Omega)<\infty$, and $\Omega$ admits hierarchical partitioning
into distinguishable subregions.
\end{axiom}

By \emph{persistent} we mean a system whose categorical identity is
preserved over time: it does not disperse, decay, or escape to spatial
infinity. The motivation is empirical: an electron bound to a nucleus,
a vibrational mode in a molecular well, a photon in a cavity, an ion
in a trap. Each is confined; in each case the unbounded idealisation
fails to persist.

\subsection{Poincar\'{e} recurrence}

\begin{theorem}[Poincar\'{e} Recurrence]
\label{thm:poincare}
Let $(\Omega,\mu,\phi_{t})$ be a measure-preserving dynamical system
with $\mu(\Omega)<\infty$. For any measurable $A\subseteq\Omega$ with
$\mu(A)>0$, almost every $x\in A$ returns to $A$ infinitely often.
\end{theorem}

\begin{proof}
Fix $T>0$ and set
$B=\{x\in A:\phi_{nT}(x)\notin A\ \forall n\geq 1\}$.
The sets $\{\phi_{nT}(B)\}_{n\geq 0}$ are pairwise disjoint: if
$\phi_{jT}(B)\cap\phi_{kT}(B)\neq\emptyset$ for some $j<k$ then some
$x\in B$ would have $\phi_{(k-j)T}(x)\in A$, contradicting the
definition of $B$. Since $\phi_{t}$ is measure-preserving,
$\mu(\phi_{nT}(B))=\mu(B)$. If $\mu(B)>0$,
$\sum_{n=0}^{\infty}\mu(\phi_{nT}(B))=\infty>\mu(\Omega)$, a
contradiction. Hence $\mu(B)=0$; almost every $x\in A$ returns. Apply
the argument iteratively to deduce infinite recurrence.
\end{proof}

\subsection{Oscillatory necessity}

\begin{theorem}[Oscillatory Necessity]
\label{thm:osc}
The only self-consistent class of dynamics on a bounded measure-finite
phase space supporting categorical observation is oscillatory.
\end{theorem}

\begin{proof}
We eliminate three alternative classes by inconsistency.

\emph{Static}. If $\phi_{t}(x)=x$ for all $t,x$, the system carries no
temporal structure; nothing distinguishes one moment from another, no
measurement can occur. This contradicts the existence of any
observation.

\emph{Monotonic}. If there exists $f:\Omega\to\R$ with $t\mapsto
f(\phi_{t}(x))$ strictly monotonic, then for any
$\varepsilon>0$ the set
$A_{\varepsilon}=\{x:f(x)<\varepsilon\}$ is exited once and never
revisited, contradicting Poincar\'{e} recurrence unless
$\mu(A_{\varepsilon})=0$ for all $\varepsilon$, which contradicts
$\mu(\Omega)<\infty$ together with $\Omega$ being measurable.

\emph{Chaotic}. If
$|\phi_{t}(x)-\phi_{t}(y)|\sim|x-y|e^{\lambda t}$ with Lyapunov
exponent $\lambda>0$, then arbitrarily close initial conditions
diverge arbitrarily in bounded time. The categorical labels
assigned by partitioning $\Omega$ are not stable: the ``same''
preparation produces different labels in finite time, violating the
reproducibility required for any observation to refer to a definite
state.

\emph{Oscillatory}. Motion that returns to neighbourhoods of previous
states (by Theorem~\ref{thm:poincare}) is bounded, recurrent, and
identity-preserving. This survives all three tests.
\end{proof}

\begin{corollary}[Discrete Mode Spectrum]
\label{cor:modes}
Bounded oscillatory dynamics admits a discrete frequency spectrum
$\{\omega_{k}\}_{k\in\N}$, the eigenvalues of the Koopman operator
$U_{t}:f\mapsto f\circ\phi_{t}$ acting on $L^{2}(\Omega,\mu)$
\citep{koopman1931}.
\end{corollary}

\begin{proof}
The Koopman operator is unitary on $L^{2}(\Omega,\mu)$ and, on a
finite-measure compact phase space, has compact resolvent and
therefore discrete spectrum. Any observable $f$ decomposes as
$f=\sum_{k}c_{k}e^{i\omega_{k}t}$ on the eigenbasis.
\end{proof}

\begin{corollary}[Bohr--Sommerfeld Quantisation]
\label{cor:bs}
The action of a fundamental oscillatory mode of frequency $\omega$ in
a bounded domain satisfies the quantisation condition
$\oint p\,dq=n\hbar\cdot 2\pi$, $n\in\N$, giving energy
$E_{n}=n\hbar\omega$. The fundamental mode has
\begin{equation}
\boxed{\;E=\hbar\omega.\;}
\label{eq:Ehbar}
\end{equation}
\end{corollary}

\begin{proof}
Single-valuedness of the wavefunction (or of the action-angle
description in the classical Bohr--Sommerfeld setting) over a closed
loop in configuration space forces the action to be an integer
multiple of $2\pi\hbar$. The fundamental mode is the $n=1$ solution.
\end{proof}

% =========================================================
\section{Partition Coordinates and Shell Capacity}
\label{sec:partition}
% =========================================================

\subsection{The four coordinates}

\begin{theorem}[Partition Coordinate System]
\label{thm:nlms}
A hierarchically nested three-dimensional oscillatory system in
bounded phase space admits a complete coordinate system
$(n,\ell,m,s)$ with
\begin{align*}
n&\in\Zp,&\ell&\in\{0,1,\ldots,n-1\},\\
m&\in\{-\ell,\ldots,+\ell\},& s&\in\{-\tfrac{1}{2},+\tfrac{1}{2}\}.
\end{align*}
\end{theorem}

\begin{proof}
We construct each coordinate from a geometric constraint on the
nested oscillatory boundary.

\emph{Step 1 (principal depth $n$).} Hierarchical nesting (Axiom
\ref{ax:bps}) provides a sequence of nested shells
$\Omega_{1}\supset\Omega_{2}\supset\cdots$. The index $n\in\Zp$
labels which shell is occupied; $n=1$ is the outermost.

\emph{Step 2 (angular complexity $\ell$).} At depth $n$ the
oscillatory boundary supports angular nodal surfaces. An angular mode
of order $\ell$ has nodal wavevector $k_{\theta}=\ell/R_{n}$ for shell
radius $R_{n}\propto n$. The mode fits the radial extent only if
$\ell\leq n-1$, giving $\ell\in\{0,1,\ldots,n-1\}$.

\emph{Step 3 (orientation $m$).} For each $\ell$, the orientation of
the angular mode in three-dimensional space is labelled by an
irreducible representation of the rotation group $SO(3)$. The
representation of angular momentum $\ell$ has dimension $2\ell+1$, with
basis labelled by $m\in\{-\ell,\ldots,+\ell\}$.

\emph{Step 4 (chirality $s$).} The fundamental group
$\pi_{1}(SO(3))=\Z_{2}$ \citep{nakahara2003} yields two topologically
distinct boundary helicities. They are labelled by half-integer
$s=\pm\tfrac{1}{2}$, reflecting the double cover $SU(2)\to SO(3)$.
\end{proof}

\subsection{The capacity formula}

\begin{theorem}[Shell Capacity]
\label{thm:cap}
The number of distinguishable partition states at depth $n$ is
\begin{equation}
C(n)=2n^{2}.
\end{equation}
\end{theorem}

\begin{proof}
At fixed $n$, sum over $\ell,m,s$:
\begin{align*}
C(n)&=\sum_{\ell=0}^{n-1}(2\ell+1)\cdot 2
=2\sum_{\ell=0}^{n-1}(2\ell+1)\\
&=2\bigl(n+2\cdot\tfrac{(n-1)n}{2}\bigr)=2n^{2},
\end{align*}
using the identity $\sum_{\ell=0}^{n-1}(2\ell+1)=n^{2}$ (the sum of
the first $n$ odd integers).
\end{proof}

\begin{corollary}[Cumulative State Count]
\label{cor:cum}
The total number of states with principal depth $\leq N$ is
\begin{equation}
\Nstate(N)=\sum_{n=1}^{N}2n^{2}=\frac{N(N+1)(2N+1)}{3}.
\label{eq:Nstate}
\end{equation}
\end{corollary}

\begin{proof}
Apply the standard sum-of-squares formula
$\sum_{n=1}^{N}n^{2}=\tfrac{N(N+1)(2N+1)}{6}$ and multiply by two.
\end{proof}

\subsection{Selection rules and exclusion}

\begin{theorem}[Selection Rules]
\label{thm:select}
Transitions between partition states satisfy
$\Delta\ell=\pm 1$, $\Delta m\in\{0,\pm 1\}$, $\Delta s=0$.
\end{theorem}

\begin{proof}
Angular complexity $\ell$ counts nodal surfaces; continuous
deformation of the oscillatory boundary creates or destroys at most
one nodal surface at a time, forcing $|\Delta\ell|=1$.
Orientation $m$ is the projection of angular momentum on a chosen
axis; continuous rotation changes it by at most one unit per
transition. Chirality is a topological invariant of the boundary
$(\pi_{1}(SO(3))=\Z_{2})$ and cannot change under continuous
deformation.
\end{proof}

\begin{theorem}[Exclusion]
\label{thm:excl}
Each partition coordinate tuple $(n,\ell,m,s)$ admits occupation
$N\in\{0,1\}$.
\end{theorem}

\begin{proof}
Two states with identical $(n,\ell,m,s)$ occupy the same partition
cell. There is no further label to distinguish a ``first'' from a
``second'' occupant; by identity of indiscernibles they are one state,
not two.
\end{proof}

% =========================================================
\section{The Partition Bijection}
\label{sec:bijection}
% =========================================================

\subsection{The map $\PhiMap$}

Let $\Parts$ denote the set of partition states with coordinates in
the ranges of Theorem~\ref{thm:nlms}.

\begin{definition}[Partition Map]
\label{def:phi}
Define $\PhiMap:\Zp\to\Parts$ as follows. Given $M\in\Zp$:
\begin{enumerate}[leftmargin=*,itemsep=0.1em]
\item Find the unique $n$ with $\Nstate(n-1)<M\leq\Nstate(n)$.
A closed-form upper bound is
$n\leq\bigl\lceil(3M/2)^{1/3}\bigr\rceil+1$.
\item Set $\delta=M-\Nstate(n-1)\in\{1,\ldots,2n^{2}\}$.
\item Enumerate the level-$n$ states in lexicographic order
$(0,0,-\tfrac{1}{2}),(0,0,+\tfrac{1}{2}),(1,-1,-\tfrac{1}{2}),\ldots$
and let $(\ell(M),m(M),s(M))$ be the $\delta$-th tuple.
\item Set $\PhiMap(M)=(n(M),\ell(M),m(M),s(M))$.
\end{enumerate}
\end{definition}

\begin{theorem}[Partition Bijection]
\label{thm:bij}
$\PhiMap:\Zp\to\Parts$ is a bijection.
\end{theorem}

\begin{proof}
\emph{Injectivity.} Suppose $\PhiMap(M_{1})=\PhiMap(M_{2})$. Then
$n(M_{1})=n(M_{2})$, so both lie in $(\Nstate(n-1),\Nstate(n)]$.
Within this block the lexicographic enumeration is strictly
monotone, so $\delta_{1}=\delta_{2}$ and therefore $M_{1}=M_{2}$.

\emph{Surjectivity.} For any $(n,\ell,m,s)\in\Parts$, set
\[
M_{0}=\Nstate(n-1)+2\sum_{\ell'=0}^{\ell-1}(2\ell'+1)+2(m+\ell)+(s+\tfrac{1}{2}).
\]
Then $\Nstate(n-1)<M_{0}\leq\Nstate(n)$, $n(M_{0})=n$, and the
level-$n$ enumeration assigns the requested tuple. Hence
$\PhiMap(M_{0})=(n,\ell,m,s)$.
\end{proof}

\begin{corollary}[Invertibility]
\label{cor:inv}
$\PhiMap^{-1}$ is given explicitly by the formula for $M_{0}$ above.
\end{corollary}

\subsection{Traversal versus existence}

\begin{definition}[Forward Sequence]
\label{def:forward}
For current partition count $M_{t}$, the \emph{forward sequence} is
$\calF(M_{t})=\{\PhiMap(M_{t}+k):k\geq 1\}\subseteq\Parts$.
\end{definition}

\begin{theorem}[Successor Existence]
\label{thm:succ}
For every $M_{t}\in\Zp$, the partition state $\PhiMap(M_{t}+1)$ is a
well-defined member of $\Parts$, independently of whether the
counting process has reached $M_{t}+1$.
\end{theorem}

\begin{proof}
By the Peano axiom on $\Zp$, $M_{t}+1\in\Zp$. Since $\PhiMap$ is
defined on all of $\Zp$ (Theorem~\ref{thm:bij}), its value at
$M_{t}+1$ is determined by Definition~\ref{def:phi}. The procedure
makes no reference to the counting process.
\end{proof}

\begin{remark}[Traversal/Existence Distinction]
\label{rem:trav}
Theorem~\ref{thm:succ} separates two predicates: (i) ``$M$ has been
traversed at time $t$'' depends on $t$; (ii) ``$\PhiMap(M)\in\Parts$''
does not. Future states are not absent from $\Parts$; they are
unvisited. This distinction is the basis for the claim that an MS
fragment is not created by the fragmentation event but is an element
of $\calF(M_{\mathrm{prec}})$ whose pre-existing partition address is
traversed.
\end{remark}

% =========================================================
\section{Time as Partition Count}
\label{sec:timecount}
% =========================================================

\subsection{Closure as a topological primitive}

\begin{definition}[Oscillatory Closure]
\label{def:closure}
An \emph{oscillatory closure} is the event of an angular coordinate
$\theta$ completing one full traversal of $[0,2\pi)$ and returning to
its initial value.
\end{definition}

\begin{proposition}[Closure is dimensionless]
\label{prop:closure}
A closure is defined without reference to duration. It requires (i) a
configuration space carrying an angular coordinate $\theta$, and (ii)
the predicate that $\theta$ has traversed $[0,2\pi)$ exactly once.
Neither (i) nor (ii) invokes any temporal unit.
\end{proposition}

\begin{proof}
Angular extent $[0,2\pi)$ is dimensionless. The predicate ``returned
to initial value'' is a property of the trajectory's geometry in
configuration space, not of how long the traversal took. Poincar\'{e}
recurrence (Theorem~\ref{thm:poincare}) guarantees that such returns
occur in any bounded oscillatory system.
\end{proof}

\subsection{Time as a derived parameter}

\begin{definition}[Time-as-Count]
\label{def:timecount}
For a bounded oscillatory system with an attached reference
oscillator at frequency $f_{\mathrm{ref}}$ stipulated by convention,
the \emph{time} of a partition count $M$ is the derived ratio
\begin{equation}
t\;:=\;\frac{M}{f_{\mathrm{ref}}}.
\label{eq:timeratio}
\end{equation}
\end{definition}

\begin{proposition}[Non-Circularity]
\label{prop:nocirc}
Equation \eqref{eq:timeratio} is not circular when read left-to-right.
$M$ counts topological closures (Definition~\ref{def:closure}), each
of which is dimensionless. The reference count $f_{\mathrm{ref}}$ is
stipulated by convention (e.g.\ $f_{\mathrm{Cs}}=9{,}192{,}631{,}770$
defining the SI second). The right-hand side is a ratio of pure
integers, and $t$ is defined to be that ratio in whichever unit the
stipulation chose.
\end{proposition}

\begin{proof}
Each closure is a topological event with no temporal content
(Proposition~\ref{prop:closure}). The count $M\in\Zp$ is an integer.
The reference count $f_{\mathrm{ref}}$ is an integer chosen by
convention. The expression $M/f_{\mathrm{ref}}$ is a rational number
of unit ``seconds'' under the chosen stipulation. No prior notion of
duration is invoked; $t$ is the named output of the ratio.
\end{proof}

\begin{theorem}[Time-Count Identity]
\label{thm:tc}
For a bounded oscillatory system at angular frequency $\omega$,
\begin{equation}
\boxed{\;\frac{dM}{dt}=\frac{\omega}{2\pi}=\frac{1}{\taup}.\;}
\label{eq:tc}
\end{equation}
\end{theorem}

\begin{proof}
One oscillation cycle traverses one partition state per period
$T_{p}=2\pi/\omega$. Hence $dM/dt=1/T_{p}=\omega/(2\pi)$. The
residence time per state is $\taup=T_{p}$, giving $1/\taup=\omega/(2\pi)$.
\end{proof}

\begin{corollary}[Strict Monotonicity]
\label{cor:mono}
$M(t_{2})>M(t_{1})$ for $t_{2}>t_{1}$.
\end{corollary}

\begin{proof}
From $dM/dt=\omega/(2\pi)>0$ in equation~\eqref{eq:tc}.
\end{proof}

\begin{corollary}[Exact Phase]
\label{cor:phase}
The oscillation phase at partition count $M$ is
$\theta(M)=2\pi M$.
\end{corollary}

\begin{proof}
$\theta=\omega t=\omega M/f_{\mathrm{ref}}$, and since
$\omega=2\pi f_{\mathrm{ref}}$ for a reference oscillator of integer
frequency, $\theta(M)=2\pi M$.
\end{proof}

% =========================================================
\section{Triple Equivalence and the GPU as Instrument}
\label{sec:triple}
% =========================================================

\subsection{Three descriptions of bounded dynamics}

\begin{theorem}[Triple Equivalence]
\label{thm:triple}
For a bounded oscillatory system with $M$ independent modes, each
partitioned into $n$ distinguishable states, the three descriptions
\begin{align*}
\text{Oscillatory}:&\quad \Omega_{\mathrm{osc}}=n^{M},\\
\text{Categorical}:&\quad \Omega_{\mathrm{cat}}=n^{M},\\
\text{Partitional}:&\quad Z=n^{M}
\end{align*}
yield the same entropy $S=\kB M\ln n$, and the conversion maps
$\mathrm{Osc}\to\mathrm{Cat}\to\mathrm{Part}\to\mathrm{Osc}$ form a
closed loop.
\end{theorem}

\begin{proof}
\emph{Osc$\to$Cat}: phase $\phi_{i}\in[0,2\pi)$ maps to category
$k_{i}=\lfloor n\phi_{i}/(2\pi)\rfloor$. Each of $M$ modes contributes
one of $n$ category labels, giving $n^{M}$ joint states.

\emph{Cat$\to$Part}: assign $\epsilon_{k}=k\epsilon_{0}$. The
single-mode partition function is
$z(\beta)=\sum_{k=0}^{n-1}e^{-\beta k\epsilon_{0}}$. In the
high-temperature limit $\kB T\gg\epsilon_{0}$, $z\to n$ and
$Z=z^{M}=n^{M}$.

\emph{Part$\to$Osc}: from $E_{k}=k\epsilon_{0}$, set
$A_{k}=\sqrt{2E_{k}/(m\omega^{2})}$ and
$x_{k}(t)=A_{k}\cos(\omega t+\phi_{0})$. The resulting phases lie in
$[0,2\pi)$ on a bin of width $2\pi/n$.

\emph{Closure}: composing the three maps returns the input phases up
to a quantisation error of at most one bin, vanishing as $n\to\infty$.

\emph{Entropy}: Boltzmann's identity $S=\kB\ln\Omega$ with
$\Omega=n^{M}$ gives $S=\kB M\ln n$ in all three descriptions.
\end{proof}

\subsection{Observation, computation, and processing}

\begin{theorem}[Observation--Computation Identity]
\label{thm:obscomp}
For a bounded oscillatory system, the following three operations are
mathematically identical:
\begin{enumerate}[leftmargin=*,itemsep=0.1em]
\item \textbf{Observation}: refining a measurement until the target
partition cell is resolved.
\item \textbf{Computation}: propagating dynamical equations until the
target partition cell is reached.
\item \textbf{Processing}: applying constraints until the target
partition cell is the unique surviving alternative.
\end{enumerate}
\end{theorem}

\begin{proof}
Each operation refines a sequence of partition selections
$\tau=(t_{1},t_{2},\ldots,t_{k})$ in $\{0,1,2\}^{k}$ (a ternary
address; see Section~\ref{sec:virtual}). The operation terminates when
the address is fully resolved at depth $k$, at which point the
partition state is unique. By Theorem~\ref{thm:triple}, the
oscillatory, categorical, and partitional realisations of the same
address are the same object; therefore the three operations operate
on the same object and terminate at the same partition state.
\end{proof}

\begin{corollary}[The Instrument is the Programme]
\label{cor:instr}
A GPU fragment shader evaluating the partition depth
$\Depth(\mathbf{x})$ at cell coordinates performs the physical
observation. It does not simulate the instrument; it is the
instrument.
\end{corollary}

\begin{proof}
A shader is a bounded oscillatory system (CPU/GPU clock at frequency
$f_{\mathrm{clk}}$) executing partition selections at rate
$dM/dt=f_{\mathrm{clk}}$. By Theorem~\ref{thm:obscomp}, the
selections it performs are operationally identical to those a
physical analyser performs on an ion in a trap. Same operation; same
output.
\end{proof}

% =========================================================
\section{The Partition Lagrangian}
\label{sec:lagrangian}
% =========================================================

\subsection{Construction from minimal assumptions}

\begin{theorem}[Partition Lagrangian]
\label{thm:Lag}
A bounded oscillatory mode with partition inertia
$\partinert=\alpha(m/z)$ in a partition-depth landscape
$\Depth(\mathbf{x},t)$ and partition-vector potential
$\Apartition(\mathbf{x})$ admits the unique Galilean-invariant
Lagrangian
\begin{equation}
\boxed{\;\Lagr_{\Depth}=\tfrac{1}{2}\partinert|\dot{\mathbf{x}}|^{2}
+\partinert\dot{\mathbf{x}}\cdot\Apartition(\mathbf{x})
-\Depth(\mathbf{x},t).\;}
\label{eq:Lagdef}
\end{equation}
\end{theorem}

\begin{proof}
We seek a Lagrangian quadratic in velocity (forced by Galilean
invariance up to total derivatives), coupling to a scalar field
$\Depth$ and a vector field $\Apartition$.

(a) The unique Galilean-invariant scalar quadratic in velocity is
$\tfrac{1}{2}\partinert|\dot{\mathbf{x}}|^{2}$.

(b) The minimal coupling of velocity to a vector potential is
$\partinert\dot{\mathbf{x}}\cdot\Apartition$.

(c) The scalar coupling is $-\Depth(\mathbf{x},t)$ with the sign
chosen so that motion descends partition depth.

There are no further first-derivative scalars; higher derivatives are
ruled out by Ostrogradsky instability. Equation~\eqref{eq:Lagdef} is
therefore unique.
\end{proof}

\begin{theorem}[Equations of Motion]
\label{thm:EOM}
The Euler--Lagrange equations of \eqref{eq:Lagdef} are
\begin{equation}
\partinert\ddot{\mathbf{x}}=-\nabla\Depth
+\partinert\dot{\mathbf{x}}\times\mathbf{B}_{\Depth},
\quad
\mathbf{B}_{\Depth}=\nabla\times\Apartition.
\label{eq:EL}
\end{equation}
\end{theorem}

\begin{proof}
Standard Euler--Lagrange algebra on \eqref{eq:Lagdef} produces
\eqref{eq:EL}; the magnetic-like term is the curl of $\Apartition$
contracted with $\dot{\mathbf{x}}$ via the vector identity
$\nabla(\dot{\mathbf{x}}\cdot\Apartition)
-d(\Apartition)/dt
=\dot{\mathbf{x}}\times(\nabla\times\Apartition)$.
\end{proof}

\begin{corollary}[Lorentz Structure as a Theorem]
\label{cor:lorentz}
With $\Depth=q\phi$ and $\Apartition=(q/\partinert)\mathbf{A}$,
\eqref{eq:EL} reduces to the Lorentz equation
$\partinert\ddot{\mathbf{x}}=q\mathbf{E}+q\dot{\mathbf{x}}\times\mathbf{B}$
with $\mathbf{E}=-\nabla\phi$, $\mathbf{B}=\nabla\times\mathbf{A}$.
The Lorentz force is the Euler--Lagrange equation for partition-depth
minimisation; no separate ``force'' assumption is required.
\end{corollary}

\subsection{The four analyser scaling laws}

We specialise $\Depth$ and $\Apartition$ to the four canonical
analyser geometries.

\textbf{Time-of-flight}: linear gradient
$\Depth_{\mathrm{TOF}}(z)=-\kappa z$, $\Apartition=\mathbf{0}$.
Solving $\partinert\ddot{z}=\kappa$ from rest over flight length $L$,
the flight time is
\begin{equation}
T_{\mathrm{TOF}}=\sqrt{\frac{2\partinert L}{\kappa}}
\;\propto\;\sqrt{m/z}.
\label{eq:tof}
\end{equation}

\textbf{Quadrupole}: oscillating saddle
$\Depth_{\mathrm{Q}}(x,y,t)=\tfrac{\kappa_{0}}{2}(x^{2}-y^{2})[U+V\cos\Omega t]$,
$\Apartition=\mathbf{0}$. Substituting $\xi=\Omega t/2$ reduces to
the Mathieu equation $\ddot u+(a-2q\cos 2\xi)u=0$ with
\begin{equation}
a=\frac{4\kappa_{0}U}{\partinert\Omega^{2}},\quad
q=\frac{2\kappa_{0}V}{\partinert\Omega^{2}}\;\propto\;\frac{1}{m/z}.
\label{eq:mathieu}
\end{equation}

\textbf{Orbitrap}: quadro-logarithmic potential
$\Depth_{\mathrm{Orb}}(r,z)=\tfrac{\kappa}{2}(z^{2}-r^{2}/2)
+\tfrac{\kappa R_{m}^{2}}{2}\ln(r/R_{m})$, $\Apartition=\mathbf{0}$.
The $z$-equation decouples:
$\partinert\ddot z=-\kappa z$, giving
\begin{equation}
\omega_{\mathrm{Orb}}=\sqrt{\kappa/\partinert}
\;\propto\;\sqrt{z/m}.
\label{eq:orb}
\end{equation}

\textbf{FT-ICR}: $\Depth=0$, vector potential
$\Apartition=(B/2)(-y,x,0)$ gives $\mathbf{B}_{\Depth}=B\hat z$. From
\eqref{eq:EL}, $\partinert\ddot{\mathbf{x}}=
\partinert\dot{\mathbf{x}}\times B\hat z$, so
\begin{equation}
\omega_{\mathrm{cyc}}=\frac{qB}{\partinert}=\frac{zeB}{m}
\;\propto\;\frac{z}{m}.
\label{eq:ftcir}
\end{equation}

\begin{theorem}[Analyser Universality]
\label{thm:univ}
Equations \eqref{eq:tof}--\eqref{eq:ftcir} are specialisations of one
Lagrangian \eqref{eq:Lagdef} under choice of partition-field topology.
\end{theorem}

\begin{proof}
Each derivation in the four paragraphs above takes \eqref{eq:Lagdef}
or \eqref{eq:EL} and a specific choice of $\Depth$ and $\Apartition$.
\end{proof}

% =========================================================
\section{Phase Coherence and Subharmonic Encoding}
\label{sec:phase}
% =========================================================

\subsection{Phase between precursor and forward sequence}

\begin{theorem}[Precursor--Fragment Phase Coherence]
\label{thm:coh}
Let a precursor ion at partition count $M_{\mathrm{prec}}$ be on a
continuous oscillatory trajectory in a partition landscape. Let
$f\in\calF(M_{\mathrm{prec}})$ be any element of the forward sequence
at count $M_{f}>M_{\mathrm{prec}}$. The oscillation phase
relationship between precursor and $f$ is
\begin{equation}
\boxed{\;\Delta\theta_{f}=2\pi\Delta M_{f}\quad\text{with}\quad
\Delta M_{f}=M_{f}-M_{\mathrm{prec}}.\;}
\label{eq:dtheta}
\end{equation}
\end{theorem}

\begin{proof}
By Corollary~\ref{cor:phase}, $\theta(M)=2\pi M$. At
$M_{\mathrm{prec}}$, $\theta_{\mathrm{prec}}=2\pi M_{\mathrm{prec}}$;
at $M_{f}$, $\theta_{f}=2\pi M_{f}$. Subtracting,
$\Delta\theta_{f}=2\pi(M_{f}-M_{\mathrm{prec}})=2\pi\Delta M_{f}$.
The argument uses only the exactness of $\theta(M)=2\pi M$ and the
fact that both states are reached on one continuous trajectory; no
assumption about the fragmentation mechanism enters.
\end{proof}

\subsection{Fragment frequencies in a fixed trap}

\begin{theorem}[Fragment Frequency Law]
\label{thm:ffreq}
In an Orbitrap trap with potential parameter $\kappa$, a fragment ion
$f$ of mass $m_{f}$ derived from a precursor of mass $m_{\mathrm{prec}}$
(same charge, same $\kappa$) oscillates at axial frequency
\begin{equation}
\omega_{f}=\omega_{\mathrm{prec}}\sqrt{\frac{m_{\mathrm{prec}}}{m_{f}}}.
\label{eq:ffreq}
\end{equation}
\end{theorem}

\begin{proof}
From \eqref{eq:orb}, $\omega=\sqrt{\kappa/\partinert}$ with
$\partinert\propto m/z$. For fixed $z$ and $\kappa$,
$\omega^{2}\propto 1/m$. Taking the ratio gives
$\omega_{f}/\omega_{\mathrm{prec}}=\sqrt{m_{\mathrm{prec}}/m_{f}}$.
\end{proof}

\subsection{Subharmonic encoding in the precursor transient}

\begin{theorem}[Fragment Subharmonic Theorem]
\label{thm:sub}
For a precursor ion with Orbitrap transient
$s_{\mathrm{prec}}(t)=A_{\mathrm{prec}}\cos(\omega_{\mathrm{prec}}t+\phi_{\mathrm{prec}})$,
each forward-sequence member $f$ contributes an additive
subharmonic component
\begin{equation}
s_{f}(t)=A_{f}\cos(\omega_{f}t+\phi_{f}),
\label{eq:sub}
\end{equation}
where $\omega_{f}=\omega_{\mathrm{prec}}\sqrt{m_{\mathrm{prec}}/m_{f}}$
(Theorem~\ref{thm:ffreq}) and
$\phi_{f}=\phi_{\mathrm{prec}}-\Delta\theta_{f}$
(Theorem~\ref{thm:coh}). Therefore the Fourier transform of
$s_{\mathrm{prec}}(t)$ contains spectral peaks at $\omega_{f}$ for
every $f$ in the forward sequence accessible during the transient.
\end{theorem}

\begin{proof}
The fragment becomes a distinct oscillator at the moment it acquires
distinguishable partition coordinates from the precursor. From that
instant, its image current is additive with the precursor's by
linearity of charge-induced detection. Its frequency is given by
Theorem~\ref{thm:ffreq}; its initial phase is the precursor's phase at
the moment of fragmentation minus the accumulated
$\Delta\theta_{f}=2\pi\Delta M_{f}$ from Theorem~\ref{thm:coh}.
Fourier-transforming the total transient yields a peak at $\omega_{f}$
with that phase.
\end{proof}

\begin{remark}[Distinguishability from Harmonics]
\label{rem:harm}
Equation \eqref{eq:ffreq} gives $\omega_{f}/\omega_{\mathrm{prec}}
=\sqrt{m_{\mathrm{prec}}/m_{f}}$, generally irrational. Pure
nonlinearity-induced harmonics of a trap appear at integer multiples
$n\omega_{\mathrm{prec}}$. The two are distinguishable in any FFT of
sufficient bin resolution.
\end{remark}

\subsection{Charge state subharmonics}

\begin{theorem}[Charge State Frequency Law]
\label{thm:zfreq}
A $z$-times-charged ion of the same neutral mass oscillates in the
Orbitrap at
\begin{equation}
\omega_{(z)}=\sqrt{z}\,\omega_{(1)}.
\label{eq:zfreq}
\end{equation}
\end{theorem}

\begin{proof}
From $\omega=\sqrt{q\kappa/m}=\sqrt{ze\kappa/m}$ at fixed neutral
mass $m$ and trap parameter $\kappa$.
\end{proof}

\begin{remark}
Charge-state and fragment subharmonics are mutually distinguishable
because they appear at $\sqrt{z}$ and $\sqrt{m_{\mathrm{prec}}/m_{f}}$
multiples respectively, with no coincidences except by accident for
small integer mass ratios.
\end{remark}

% =========================================================
\section{Virtual Substates and Mean Recovery}
\label{sec:virtual}
% =========================================================

\subsection{S-entropy coordinate space}

\begin{definition}[S-Entropy Coordinates]
\label{def:Sspace}
The \emph{S-entropy coordinate space} is
$\Sspace=[0,1]^{3}$ with coordinates $(S_{k},S_{t},S_{e})$. We
identify three orthogonal physical projections:
\begin{itemize}[nosep,leftmargin=*]
\item $S_{k}$: knowledge (mass/charge scale).
\item $S_{t}$: temporal (retention or arrival time scale).
\item $S_{e}$: evolution (fragmentation state).
\end{itemize}
\end{definition}

\begin{definition}[Ternary Address]
\label{def:ternary}
A \emph{ternary address} of depth $k$ is a sequence
$\tau=(t_{1},\ldots,t_{k})\in\{0,1,2\}^{k}$. The set of all
depth-$k$ addresses has cardinality $3^{k}$.
\end{definition}

\begin{definition}[Trit-Axis Mapping]
\label{def:trit}
The trit at position $i$ refines axis $S_{k},S_{t},S_{e}$ according
to $i\bmod 3$. Cell $C_{\tau}$ at depth $k$ has volume $3^{-k}$ in
$\Sspace$ and centre
\begin{equation}
\mathbf{r}(\tau)=\Bigl(r_{k}(\tau),r_{t}(\tau),r_{e}(\tau)\Bigr),
\label{eq:rcell}
\end{equation}
where each $r_{a}(\tau)=\sum_{i\in I_{a}}t_{i}\cdot 3^{-\lceil
i/3\rceil}\cdot \tfrac{1}{2}+\tfrac{1}{2}\cdot 3^{-\lceil i/3\rceil}$
with $I_{a}$ the trit positions assigned to axis $a$.
\end{definition}

\begin{theorem}[Position--Trajectory Duality]
\label{thm:dual}
A ternary address $\tau$ simultaneously encodes (i) the cell
$C_{\tau}$ (final position) and (ii) the sequence of refinements
(trajectory). The map is bijective.
\end{theorem}

\begin{proof}
Reading $\tau$ left-to-right gives the sequence of partition
selections starting from $[0,1]^{3}$. The cell $C_{\tau}$ is the
intersection of those selections. Different addresses yield
different cells (injectivity from disjointness of subdivisions); every
cell at depth $k$ has a unique address (surjectivity from exhaustive
subdivision).
\end{proof}

\subsection{Virtual substates}

\begin{definition}[Virtual Substate]
\label{def:virtual}
For a partition value $s\in[0,1]$ and integer $N\geq 2$, a
\emph{virtual substate} of $s$ is an ordered tuple
$(s_{1},\ldots,s_{N})\in\R^{N}$ satisfying the \emph{mean-recovery
constraint}
\begin{equation}
\bar s:=\frac{1}{N}\sum_{j=1}^{N}s_{j}=s.
\label{eq:meanrec}
\end{equation}
A virtual substate is \emph{physical} if every $s_{j}\in[0,1]$ and
\emph{off-shell} otherwise.
\end{definition}

\begin{theorem}[Existence of Off-Shell Decompositions]
\label{thm:exoff}
For any $s\in(0,1)$ and $N\geq 2$, the set of off-shell virtual
substates of $s$ has positive Lebesgue measure on the constraint
hyperplane $\{(s_{1},\ldots,s_{N}):\sum s_{j}=Ns\}$.
\end{theorem}

\begin{proof}
The constraint defines an affine $(N-1)$-dimensional hyperplane in
$\R^{N}$. The set of physical substates is the bounded polytope
$\{[0,1]^{N}\cap\text{constraint}\}$, of finite Lebesgue measure. The
complement on the constraint hyperplane has infinite measure for any
$s\in(0,1)$, so off-shell substates dominate.
\end{proof}

\begin{theorem}[Miracle Principle]
\label{thm:mir}
For any $s\in[0,1]$, $N\geq 2$, and any choice of
$(s_{1},\ldots,s_{N-1})\in\R^{N-1}$, the value
$s_{N}=Ns-\sum_{j=1}^{N-1}s_{j}$ is uniquely determined and the
mean-recovery constraint holds.
\end{theorem}

\begin{proof}
Solve \eqref{eq:meanrec} for $s_{N}$.
\end{proof}

\subsection{Unconstrained subtask theorem}

\begin{theorem}[Unconstrained Subtask]
\label{thm:unsub}
Let $\xi$ be an expression evaluating to global value $g\in[0,1]$ in
a fixed evaluation calculus, and let $\eta$ be any subexpression of
$\xi$. The substitution
$\xi'=C[\eta\mapsto\eta_{1}+\eta_{2}-\eta_{1}]$ where $\eta_{1}$ has
arbitrary local value $\sigma\in\R$ and $\eta_{2}$ has
$\mathrm{value}(\eta_{2})=\mathrm{value}(\eta)$, satisfies
$\mathrm{value}(\xi')=g$ regardless of $\sigma$.
\end{theorem}

\begin{proof}
The arithmetic identity $\eta_{1}+\eta_{2}-\eta_{1}=\eta_{2}$ holds
for any $\eta_{1}$. The surrounding context $C[\cdot]$ is
unchanged. By compositionality of evaluation, $\xi'$ evaluates to the
same value as the original $\xi$.
\end{proof}

\begin{corollary}[Free Local Virtual Components]
\label{cor:freeloc}
Under preserved global mean-recovery, individual virtual components
$s_{j}$ are unconstrained.
\end{corollary}

\subsection{Connection to off-shell propagators}

\begin{remark}[Formal Identification with Feynman Propagators]
\label{rem:feyn}
In quantum field theory the Feynman propagator
$D_{F}(q)=i/(q^{2}-m^{2}c^{4}+i\varepsilon)$ admits virtual
4-momenta $q$ with $q^{2}\neq m^{2}c^{4}$ (off-shell), subject to
energy--momentum conservation $\sum q_{i}=Q$ at each vertex
\citep{peskin1995}. The mean-recovery constraint $\bar s=s$ on
virtual substates is the same statement: off-shell intermediate
values are free, but their (suitably weighted) sum equals the
physical external value. The identification is formal, not a
metaphor: both are intermediate decompositions outside the physical
on-shell domain whose aggregate satisfies a physical constraint.
\end{remark}

% =========================================================
\section{The Stacked Virtual Tensor}
\label{sec:tensor}
% =========================================================

\subsection{Four orthogonal decomposition dimensions}

An MS ion admits four natural dimensions of virtual decomposition,
each carrying an independent mean-recovery constraint.

\textbf{D1 (instrument basis)}: same ion, four trap geometries,
indexed $i\in\{1,2,3,4\}$ for Orbitrap, FT-ICR, TOF, quadrupole.

\textbf{D2 (charge state)}: $z\in\{1,\ldots,Z\}$ for the same neutral
mass.

\textbf{D3 (polarity)}: $k\in\{0,1\}$ for $\pm$ ion mode (chirality
$s=\pm\tfrac{1}{2}$).

\textbf{D4 (time)}: $\ell\in\{1,\ldots,n_{t}\}$ for temporal substates
relative to the current measurement.

\begin{definition}[Virtual Partition Tensor]
\label{def:tensor}
The \emph{virtual partition tensor} is the rank-4 array
$V_{ijkl}$ with $V_{ijkl}\in\R^{3}$, the virtual partition
coordinate at the $(i,j,k,\ell)$ corner of the stacked decomposition.
\end{definition}

\begin{theorem}[Tensor Mean-Recovery]
\label{thm:tensor}
The physical ion's partition coordinate
$\mathbf v_{\mathrm{phys}}\in[0,1]^{3}$ satisfies
\begin{equation}
\mathbf v_{\mathrm{phys}}=\frac{1}{N}\sum_{i,j,k,\ell}V_{ijkl},
\quad N=4\cdot Z\cdot 2\cdot n_{t}.
\label{eq:tensormean}
\end{equation}
\end{theorem}

\begin{proof}
Theorem~\ref{thm:exoff} and Definition~\ref{def:virtual} extend
coordinate-wise to the four-index sum.
\end{proof}

\subsection{Degrees of freedom and information content}

\begin{proposition}[Degrees of Freedom]
\label{prop:dof}
The tensor mean-recovery constraint is one vector equation in
$\R^{3}$, i.e.\ three scalar constraints. The number of free
parameters is
\begin{equation}
\mathrm{dof}(V)=3(N-1).
\label{eq:dof}
\end{equation}
\end{proposition}

\begin{proof}
$V$ has $3N$ scalar entries; three are removed by \eqref{eq:tensormean}.
\end{proof}

\begin{example}[Peptide Scale]
For $Z=500$, $n_{t}=100$:
$N=4\cdot 500\cdot 2\cdot 100=4\times 10^{5}$, giving
$\mathrm{dof}(V)=3(4\times 10^{5}-1)\approx 1.2\times 10^{6}$ free
parameters. These encode every entry of the four decomposition
dimensions: charge envelope, fragment substates, polarity complement,
ionisation history. They are present in the single-precursor measurement.
\end{example}

% =========================================================
\section{Backward Trajectory Completion}
\label{sec:backward}
% =========================================================

\subsection{Hierarchy and complexity}

\begin{definition}[Ternary Refinement Hierarchy]
\label{def:hier}
A \emph{ternary refinement hierarchy} of depth $k$ in $\Sspace$ is a
sequence of partitions $\{P_{0},\ldots,P_{k}\}$ with
$|P_{j+1}|=3|P_{j}|$, $|P_{0}|=1$, and each block of $P_{j}$ refined
into three children. Total leaves: $N=3^{k}$.
\end{definition}

\begin{theorem}[Discrete Lower Bound]
\label{thm:lower}
Any algorithm solving the trajectory-completion problem using only
equality and membership queries on $\Sspace$ requires $\Omega(N)$
queries in the worst case.
\end{theorem}

\begin{proof}
An adversary can withhold partition structure until forced to reveal
it. Distinguishing the leaf block of every other leaf requires
$\Omega(N\log_{2}3)$ bits of information, each query yielding at most
$O(1)$ bits. The bound follows.
\end{proof}

\subsection{Logarithmic completion with virtual substates}

\begin{theorem}[Backward Navigation Speedup]
\label{thm:logn}
Backward trajectory completion using virtual substates terminates in
exactly $\log_{3}N=k$ steps.
\end{theorem}

\begin{proof}
Each step computes the parent cell of the current address by
truncating the last trit of the base-3 expansion of $(S_{k},S_{t},S_{e})$.
This is $\mathcal{O}(1)$ per step. Starting at depth $k$ and ending
at depth $0$, exactly $k$ steps are required.
\end{proof}

\begin{theorem}[Collapse Without Virtual Substates]
\label{thm:collapse}
Restricting decompositions to physical (on-shell) sub-coordinates
collapses backward completion to $\Theta(N)$.
\end{theorem}

\begin{proof}
A physical-only step at depth $j$ requires geodesic resolution in the
continuous metric finer than the cell width at depth $j-1$, which is
attainable only via decompositions whose components escape $[0,1]^{3}$
(Theorem~\ref{thm:exoff}). Without virtual components, the parent
must be located by enumerating children of candidate parents, summing
to $\sum_{j=0}^{k}3^{j}=\Theta(N)$.
\end{proof}

\begin{algorithm}[!htbp]
\caption{Backward Completion (virtual)}
\label{alg:back}
\begin{algorithmic}[1]
\Function{Complete}{addr $\tau\in\{0,1,2\}^{k}$}
\State trajectory $\mathcal{T}\gets[\tau]$
\State $j\gets k$
\While{$j>0$}
\State $\tau\gets\tau[1{:}j-1]$
\Comment{drop last trit (parent)}
\State $\mathcal{T}.\mathrm{prepend}(\tau)$
\State $j\gets j-1$
\EndWhile
\State \Return $\mathcal{T}$
\EndFunction
\end{algorithmic}
\end{algorithm}

\begin{theorem}[Path Opacity]
\label{thm:path}
Let $\mathcal{T}_{1},\mathcal{T}_{2}$ be two backward trajectories
sharing endpoint $\tau_{k}$ but differing in intermediate virtual
decompositions. Any metric invariant computable at $\tau_{k}$
(coordinate, distance to origin, Riemannian volume) is equal on both.
\end{theorem}

\begin{proof}
Endpoint-defined invariants are functions of $\tau_{k}$ alone. By
Theorem~\ref{thm:unsub}, intermediate virtual values can vary
freely under preserved endpoint, so two trajectories with the same
$\tau_{k}$ yield the same invariants.
\end{proof}

% =========================================================
\section{The Six-Pass Shader Pipeline}
\label{sec:shader}
% =========================================================

\subsection{Mapping the ion journey to shader passes}

By Corollary~\ref{cor:instr}, a GPU shader evaluating partition depth
at cell coordinates performs the physical observation. We map each
stage of the ion journey to one shader pass.

\begin{table}[H]
\centering
\caption{Six-pass pipeline.}
\label{tab:pipeline}
\small
\begin{tabular}{lll}
\toprule
\textbf{Pass} & \textbf{Stage} & \textbf{Operation} \\
\midrule
1 & Ionisation & Partition creation $\Delta\Depth$\\
2 & Ion optics & Gradient descent $-\nabla\Depth$\\
3 & Mass analysis & Topology specialisation\\
4 & Fragmentation & $\Depth$ redistribution\\
5 & Detection & Partition completion\\
6 & Signal & Categorical state flux\\
\bottomrule
\end{tabular}
\end{table}

\subsection{Pass specifications}

\textbf{Pass 1 (Ionisation)}. Input: neutral molecule parameters.
Operation: evaluate $\Delta\Depth=\kB T\ln b\cdot\sum_{i}\log_{b}k_{i}$
for creating $z$ unfilled partition slots. Ionisation energy is
$\mathrm{IE}=T\kB\ln b\cdot\Delta\Depth$. Output: ion parameters
$(m/z,z)$.

\textbf{Pass 2 (Ion optics)}. Input: ion at source exit. Operation:
evaluate $\Depth(\mathbf{x})$ from electrode geometry and integrate
$\dot{\mathbf{x}}=-\gamma\nabla\Depth$. Output: ion position at
analyser entrance.

\textbf{Pass 3 (Mass analysis)}. Input: ion at analyser entrance.
Operation: evaluate analyser-specific $\Depth(\mathbf{x},t)$ from
\eqref{eq:tof}--\eqref{eq:ftcir}. Output: analyser observable.

\textbf{Pass 4 (Fragmentation)}. Input: parent ion at $\Depth_{p}$,
collision energy $E_{c}$. Operation: redistribute $\Depth_{p}$ to
$\{\Depth_{f_{i}}\}$ subject to $\sum_{i}\Depth_{f_{i}}=\Depth_{p}$ and
selection rules (Theorem~\ref{thm:select}). Output: fragment list.

\textbf{Pass 5 (Detection)}. Input: ion at detector with $\Depth$.
Operation: compute $\Delta\Depth_{\mathrm{compl}}=\Depth-\Depth_{\min}$;
signal amplitude $\propto\Delta\Depth_{\mathrm{compl}}$. Output: signal
amplitude.

\textbf{Pass 6 (Signal)}. Input: detection event, network topology.
Operation: propagate categorical state through phase-locked
electronics at $v\sim\sqrt{G/\rho_{m}}$. Output: digitised signal.

\begin{theorem}[Pipeline Completeness]
\label{thm:complete}
The six-pass pipeline reproduces every observable a physical
mass spectrometer produces for an ion in $\Parts$, with zero forces
and zero adjustable parameters beyond the trap geometry constants.
\end{theorem}

\begin{proof}
Pass 1 produces correct $\mathrm{IE}$ by depth-energy equivalence.
Pass 2 produces correct trajectories from gradient flow of $\Depth$.
Pass 3 produces the correct scaling \eqref{eq:tof}--\eqref{eq:ftcir}.
Pass 4 conserves $\Depth$ and respects selection rules. Pass 5
produces signals proportional to $\Delta\Depth$. Pass 6 produces
realistic noise via partition lag. Every stage is a derivation from
\eqref{eq:Lagdef}; no external assumption enters.
\end{proof}

% =========================================================
\section{The Five Falsifiable Experiments}
\label{sec:experiments}
% =========================================================

The validation programme falsifies five claims of the framework.
Each claim is a theorem proved above; each experiment is a numerical
operation on real \texttt{mzML} or transient data with a predicted
statistic, a pass criterion at the framework's predicted precision,
and a clear failure mode.

\subsection{E1: Bijection and capacity}

\begin{experiment}[Partition Bijection on MS Data]
\label{exp:e1}
\textbf{Theorem under test:} Theorems~\ref{thm:cap},
\ref{thm:bij}, and Corollary~\ref{cor:cum}.\\
\textbf{Input:} \texttt{mzML} with $N_{\mathrm{ion}}$ detected ions.\\
\textbf{Procedure:}
\begin{enumerate}[nosep,leftmargin=*]
\item For each ion with observed $m/z$, compute
$M_{i}=\lfloor f_{\mathrm{ref}}\cdot t_{\mathrm{flight},i}\rfloor$
when raw timing is available, otherwise
$n_{i}=\lfloor\sqrt{(m/z)_{i}/m_{\mathrm{ref}}}\rfloor+1$,
$\delta_{i}=\lceil(2n_{i}^{2})\cdot S_{e,i}\rceil$, and
$M_{i}=\Nstate(n_{i}-1)+\delta_{i}$.
\item Apply $\PhiMap$ to obtain $(n_{i},\ell_{i},m_{i},s_{i})$.
\item Check capacity $\delta_{i}\leq 2n_{i}^{2}$ and selection
$0\leq\ell_{i}<n_{i}$, $|m_{i}|\leq\ell_{i}$,
$s_{i}\in\{\pm\tfrac{1}{2}\}$ for all ions.
\item Check injectivity: $M_{i}=M_{j}\Leftrightarrow i=j$.
\end{enumerate}
\textbf{Pass criterion:} 100\% of ions satisfy capacity, selection,
and injectivity.\\
\textbf{Failure mode:} any ion outside the bounds falsifies
Theorem~\ref{thm:cap}; any collision falsifies Theorem~\ref{thm:bij}.
\end{experiment}

\subsection{E2: Time-count identity}

\begin{experiment}[Time-Count Identity in a Trap]
\label{exp:e2}
\textbf{Theorem under test:} Theorem~\ref{thm:tc}.\\
\textbf{Input:} Orbitrap transient with stipulated reference
oscillator frequency $f_{\mathrm{ref}}$, $N_{p}$ peaks at frequencies
$\{\omega_{k}\}$.\\
\textbf{Procedure:}
\begin{enumerate}[nosep,leftmargin=*]
\item For each peak compute $M_{k}(T)=f_{k}\cdot T$ over the transient
duration $T$, where $f_{k}=\omega_{k}/(2\pi)$.
\item Compute the residual $r_{k}=M_{k}(T)-\omega_{k}T/(2\pi)$ in
machine precision.
\end{enumerate}
\textbf{Pass criterion:} $\max_{k}|r_{k}|<10^{-12}$.\\
\textbf{Failure mode:} any $|r_{k}|$ above floating-point noise
falsifies the identity.
\end{experiment}

\subsection{E3: Analyser scaling laws}

\begin{experiment}[Four Analyser Scaling Laws]
\label{exp:e3}
\textbf{Theorem under test:} Theorem~\ref{thm:univ} and
equations~\eqref{eq:tof}--\eqref{eq:ftcir}.\\
\textbf{Input:} reference dataset of $N_{c}$ compounds with measured
flight times (TOF), Orbitrap frequencies (Orb), Mathieu stability
($a,q$) (Quad), and cyclotron frequencies (FT-ICR).\\
\textbf{Procedure:} for each analyser type, fit the predicted
scaling and compute the maximum residual:
\begin{align*}
\mathrm{res}_{\mathrm{TOF}}&=\max_{c}\bigl|T_{c}-K_{\mathrm{T}}\sqrt{(m/z)_{c}}\bigr|/T_{c},\\
\mathrm{res}_{\mathrm{Orb}}&=\max_{c}\bigl|\omega_{c}-K_{\mathrm{O}}\sqrt{(z/m)_{c}}\bigr|/\omega_{c},\\
\mathrm{res}_{\mathrm{ICR}}&=\max_{c}\bigl|\omega_{c}-K_{\mathrm{I}}(z/m)_{c}\bigr|/\omega_{c},\\
\mathrm{res}_{\mathrm{Q}}&=\max_{c}\bigl|q_{c}-K_{\mathrm{Q}}(m/z)_{c}^{-1}\bigr|/q_{c}.
\end{align*}
\textbf{Pass criterion:} all four residuals $<10^{-4}$ on a dataset
of $N_{c}\geq 30$ compounds spanning at least one decade in $m/z$.\\
\textbf{Failure mode:} any residual above $10^{-4}$ falsifies
\eqref{eq:Lagdef} for the given analyser.
\end{experiment}

\subsection{E4: Subharmonic encoding}

\begin{experiment}[Single-Transient Fragment Encoding]
\label{exp:e4}
\textbf{Theorem under test:} Theorems~\ref{thm:coh},
\ref{thm:ffreq}, \ref{thm:sub}.\\
\textbf{Input:} raw Orbitrap transient $s(t)$ of one precursor at
$\omega_{\mathrm{prec}}$, paired with the corresponding centroided
MS\textsuperscript{2} fragment list
$\{(m_{f_{i}},I_{f_{i}})\}_{i=1}^{N_{f}}$.\\
\textbf{Procedure:}
\begin{enumerate}[nosep,leftmargin=*]
\item Compute predicted subharmonic frequencies
$\omega_{f_{i}}^{\mathrm{pred}}=\omega_{\mathrm{prec}}\sqrt{m_{\mathrm{prec}}/m_{f_{i}}}$
and predicted phases
$\phi_{f_{i}}^{\mathrm{pred}}=\phi_{\mathrm{prec}}-2\pi\Delta M_{f_{i}}$.
\item Compute $\hat s(\omega)$ by FFT with bin resolution
$\Delta\omega<\min_{i}|\omega_{f_{i}}-\omega_{\mathrm{prec}}|/10$.
\item For each predicted $\omega_{f_{i}}^{\mathrm{pred}}$, search
$\hat s$ within $\pm 3\Delta\omega$ for a peak above the local noise
floor at threshold SNR $\geq 5$.
\item Record measured peak frequency and phase
$(\omega_{f_{i}}^{\mathrm{meas}},\phi_{f_{i}}^{\mathrm{meas}})$.
\item Compute frequency residual
$r_{i}^{\omega}=|\omega_{f_{i}}^{\mathrm{meas}}-\omega_{f_{i}}^{\mathrm{pred}}|/\omega_{f_{i}}^{\mathrm{pred}}$
and phase residual
$r_{i}^{\phi}=|\phi_{f_{i}}^{\mathrm{meas}}-\phi_{f_{i}}^{\mathrm{pred}}|\bmod\pi$.
\end{enumerate}
\textbf{Pass criterion:} for $\geq 80\%$ of MS\textsuperscript{2}
fragments above the SNR floor,
$r_{i}^{\omega}<10^{-5}$ and $r_{i}^{\phi}<0.1$~rad.\\
\textbf{Failure mode:} subharmonics absent at predicted positions,
or present but with phases inconsistent with $2\pi\Delta M$, or
present at wrong frequencies. Either falsifies one of
Theorems~\ref{thm:coh}, \ref{thm:ffreq}, \ref{thm:sub}.
\end{experiment}

\begin{remark}[Why this is the headline]
E4 is the experiment that distinguishes the partition picture from
the force picture observationally. The two pictures agree on the
$m/z$ scaling laws; they disagree on whether fragment information is
present in a single MS\textsuperscript{1} transient. If E4 passes,
no further evidence is needed: a single transient is operationally
equivalent to an MS\textsuperscript{2} measurement.
\end{remark}

\subsection{E5: Backward completion complexity}

\begin{experiment}[Logarithmic Backward Completion]
\label{exp:e5}
\textbf{Theorem under test:} Theorems~\ref{thm:logn},
\ref{thm:collapse}, \ref{thm:path}.\\
\textbf{Input:} for $N_{c}$ ions with annotated S-coordinates
$(S_{k},S_{t},S_{e})$, the corresponding ternary addresses at depth
$k=18$ (six trits per axis).\\
\textbf{Procedure (virtual path):}
\begin{enumerate}[nosep,leftmargin=*]
\item For each ion, apply Algorithm~\ref{alg:back} to obtain the
backward trajectory and record step count $K_{i}^{\mathrm{virt}}$.
\item For an opacity check, generate two distinct random
intermediate virtual decompositions sharing the endpoint and
record metric invariants at the endpoint.
\end{enumerate}
\textbf{Procedure (physical-only path):}
\begin{enumerate}[nosep,leftmargin=*]
\item Repeat backward completion restricted to physical
substates (mean-recovery with all components in $[0,1]$).
\item Record step count $K_{i}^{\mathrm{phys}}$.
\end{enumerate}
\textbf{Pass criterion:}
$K_{i}^{\mathrm{virt}}=k$ exactly for all $i$;
$\langle K_{i}^{\mathrm{phys}}/(3^{k})\rangle$ within a constant
factor of unity; endpoint invariants in the opacity check identical
for both random intermediate decompositions to floating-point
precision.\\
\textbf{Failure mode:} virtual completion exceeding $k$ steps
falsifies Theorem~\ref{thm:logn}; physical-only completion in
sub-linear steps falsifies Theorem~\ref{thm:collapse}; invariants
differing between paths falsifies Theorem~\ref{thm:path}.
\end{experiment}

% =========================================================
\section{Data Contracts, JSON Records, and Pass Criteria}
\label{sec:contracts}
% =========================================================

\subsection{Common record schema}

Each experiment persists a JSON record of the form
\begin{verbatim}
{
  "experiment": "E1|E2|E3|E4|E5",
  "theorem_id": <int>,
  "input_dataset": <str>,
  "n_samples": <int>,
  "predicted": { ... },
  "measured": { ... },
  "residuals": { "max": <float>,
                 "rms": <float> },
  "pass": <bool>,
  "elapsed_seconds": <float>,
  "seed": <int>
}
\end{verbatim}

\subsection{Per-experiment specifications}

\textbf{E1.} \emph{Inputs}: \texttt{mzML} with peaks. \emph{Predicted}:
$M_{i}\in\Zp$, $(n,\ell,m,s)$. \emph{Measured}: same, by
$\PhiMap$. \emph{Pass}: 100\% of $N_{\mathrm{ion}}$ satisfy capacity
and selection; no collisions.

\textbf{E2.} \emph{Inputs}: transient peak list, $f_{\mathrm{ref}}$,
$T$. \emph{Predicted}: $M_{k}^{\mathrm{pred}}=f_{k}T$. \emph{Measured}:
$M_{k}^{\mathrm{meas}}=\omega_{k}T/(2\pi)$. \emph{Pass}:
$\max|r_{k}|<10^{-12}$.

\textbf{E3.} \emph{Inputs}: $\geq 30$ compounds with all four
analyser observables. \emph{Predicted}: scaling \eqref{eq:tof}--\eqref{eq:ftcir}.
\emph{Measured}: instrument data. \emph{Pass}: all residuals $<10^{-4}$.

\textbf{E4.} \emph{Inputs}: raw Orbitrap transient and MS\textsuperscript{2}
fragment list. \emph{Predicted}: $(\omega_{f_{i}}^{\mathrm{pred}},
\phi_{f_{i}}^{\mathrm{pred}})$. \emph{Measured}: FFT-extracted peaks
within SNR$\geq 5$. \emph{Pass}: $\geq 80\%$ fragments matched with
$r^{\omega}<10^{-5}$ and $r^{\phi}<0.1$~rad.

\textbf{E5.} \emph{Inputs}: $N_{c}$ ternary addresses at depth $k=18$.
\emph{Predicted}: $K^{\mathrm{virt}}=k$,
$K^{\mathrm{phys}}\sim 3^{k}$. \emph{Measured}: step counts.
\emph{Pass}: identity for virtual; linear-in-$N$ scaling for physical;
opacity invariants equal.

\subsection{Test grids and dataset choice}

Recommended datasets:
\begin{itemize}[nosep,leftmargin=*]
\item E1: HMDB-derived \texttt{mzML} ($\sim 2{,}847$ compounds);
LIPID MAPS ($\sim 892$ species); MassBank ($\sim 532$ compounds).
\item E2: any Orbitrap transient with manufacturer-reported reference
oscillator $f_{\mathrm{ref}}$.
\item E3: 39 NIST compounds spanning all four analyser types.
\item E4: a paired raw transient and MS\textsuperscript{2} scan on
the same precursor; HCD or CID fragmentation.
\item E5: any subset of E1's compounds with depth-18 addresses.
\end{itemize}

\subsection{Suite-level acceptance}

The full validation suite is accepted if and only if all five
experiments report \texttt{PASS}. A single \texttt{FAIL} falsifies
one of the named theorems and localises the failure precisely: E1
(bijection), E2 (time-count), E3 (Lagrangian), E4 (phase coherence
and subharmonics), E5 (backward completion).

\subsection{Reproducibility}

All experiments are deterministic given the random seed. JSON records
include the seed and elapsed time. A separate master record
$\texttt{master\_results.json}$ aggregates per-experiment pass flags
and the wall-clock time of the whole suite.

% =========================================================
\section{Discussion}
\label{sec:discussion}
% =========================================================

\subsection{What this blueprint proves and what it does not}

The derivations in Sections~\ref{sec:axiom}--\ref{sec:shader}
establish that, granted the single Bounded Phase Space Axiom
(Axiom~\ref{ax:bps}), the partition-depth picture of a mass
spectrometer is internally consistent, dimensionally honest, and
free of circularity (Proposition~\ref{prop:nocirc}). They do not
establish that the picture is physically correct; the experiments
of Section~\ref{sec:experiments} are designed to falsify it if it is
wrong.

\subsection{Scope and limitations}

\begin{enumerate}[nosep,leftmargin=*]
\item \textbf{Non-integrable systems.} The Lagrangian derivation
(Theorem~\ref{thm:Lag}) assumes integrability in the sense that
action-angle variables exist. KAM-type robustness arguments are
necessary to extend to near-integrable cases.
\item \textbf{Effective-charge approximations.} The atomic
ionisation-energy comparisons used as historical motivation for
$C(n)=2n^{2}$ involve screening corrections we do not model here.
\item \textbf{Relativistic regime.} The analyser scaling derivations
assume $v\ll c$. Relativistic generalisation requires the
Lorentzian metric in place of the Euclidean kinetic term.
\item \textbf{Detector noise.} Pass 6 of the shader pipeline
specifies the categorical-state-flux noise model. Practical detection
electronics may introduce additional non-categorical noise (thermal,
shot) that broadens FFT peaks in E4.
\end{enumerate}

\subsection{Connection to the wider categorical computing framework}

The partition Lagrangian, time-count identity, and triple equivalence
extend beyond mass spectrometry. Any bounded oscillatory measurement
(NMR, X-ray crystallography, IR/Raman, FT-ICR mass spectrometry,
Orbitrap, ion mobility) admits the same six-pass mapping, with the
specific $\Depth$-topology determining the analyser type. The
blueprint here is implemented for mass spectrometry but is portable
to other domains by changing the $\Depth$-field uniforms in Pass 3.

\subsection{What success and failure mean}

If E1--E5 all pass, the partition picture is operationally validated
on real data: mass spectrometry is determined by partition addresses,
single transients carry MS\textsuperscript{2} information, backward
completion is logarithmic, and the six-pass shader pipeline is a
complete instrument.

If E4 fails while E1--E3, E5 pass, the partition framework remains
consistent but the precise phase-coherence claim is empirically
wrong: the time-domain transient does not carry fragment information.
This would be a precise, localisable falsification: not of the whole
framework, but of Theorem~\ref{thm:coh} or
Theorem~\ref{thm:sub}.

If E1 fails on real data---for instance, a violation of the capacity
bound $\delta_{i}\leq 2n_{i}^{2}$---the partition coordinate system
of Theorem~\ref{thm:nlms} would be falsified, and all downstream
results would require revision.

\subsection{Outlook}

The blueprint as specified can be implemented in Python with
\texttt{numpy} and \texttt{scipy} for E1--E3, E5, and with
\texttt{numpy.fft} plus a Thermo \texttt{.raw} parser for E4. The
target is a single \texttt{python -m validate\_all} entry point
producing a JSON suite analogous to the 14-experiment artefact in
the categorical computing framework. The full implementation is left
to the companion software paper; the present paper specifies what is
to be implemented and what counts as success.

% =========================================================
\section*{Acknowledgements}
% =========================================================

This work is part of the bounded-phase-space programme at the
TU Munich Chair of Proteomics and Bioanalytics. No external funding
was used. All datasets cited are publicly available.

% =========================================================
\bibliographystyle{plainnat}
\bibliography{references}
% =========================================================

\end{document}
